% Κοινή κατασκευή AttentionBox / ThinkBox / HistoryBox.
% Οι δηλώσεις στο τέλος αλλάζουν φόντο, εικονίδιο και ετικέτα.
% Ο κοινός BookCalloutAccent χρωματίζει τίτλους, εικονίδια και διακόσμηση.
% Το ThinkBox χρησιμοποιεί αγκύλες χωρίς γέμισμα· τα άλλα δύο πλαίσιο.
% Γεωμετρία, τίτλος, διακόσμηση και συμπεριφορά
% αλλαγής σελίδας υλοποιούνται μία φορά για όλα τα κουτιά.
% Προαπαιτούμενα: BookCallout, BookCalloutCenteredIcon, BookMediumLabelFont
% και οι τρεις σχεδιαστικές εντολές εικονιδίων του universalbook.sty.

\providecommand{\BookCalloutLabelSize}{9.2}
\providecommand{\BookCalloutLabelLeading}{11.5}

% Ο τίτλος είναι σώμα του κουτιού (όχι δεύτερο ανεξάρτητο title box).
% Το προαιρετικό θέμα διατηρεί τη σύνταξη HistoryBox[Θέμα].
\newcommand{\BookCalloutHeading}[3]{%
  \noindent{\BookMediumLabelFont
    \fontsize{\BookCalloutLabelSize}{\BookCalloutLabelLeading}\selectfont
    \color{#1}#2\ifstrempty{#3}{}{ — #3}}\par\smallskip\small
}

% Πέντε παράμετροι: φόντο / έμφαση / εικονίδιο / τίτλος / θέμα.
% Το ίδιο overlay χρησιμοποιείται σε ενιαίο και πρώτο σπασμένο τμήμα.
% Τα επόμενα τμήματα κρατούν τη γεωμετρία χωρίς να επαναλαμβάνουν icon.
\tcbset{BookCalloutRender/.style n args={5}{%
  BookCallout,
  colback=#1,
  borderline={\BookBoxThinRuleWidth}{0pt}{#2},
  borderline west={\BookCalloutAccentRuleWidth}{0pt}{#2},
  before upper={\BookCalloutHeading{#2}{#4}{#5}},
  overlay unbroken={\BookCalloutCenteredIcon{#3}},
  overlay first={\BookCalloutCenteredIcon{#3}}
}}

% Δύο διανυσματικές αγκύλες περιέχουν και τη ζώνη εικονιδίου και το κείμενο.
% Παράμετροι: χρώμα / επάνω εσοχή / κάτω εσοχή. Μηδενική εσοχή σημαίνει
% ανοιχτή συνέχεια στην αλλαγή σελίδας, χωρίς οριζόντιο άκρο αγκύλης.
\newcommand{\BookCalloutBrackets}[3]{%
  \coordinate (BookBracketNW) at
    ([xshift=\BookThinkBracketInset,yshift=-#2]frame.north west);
  \coordinate (BookBracketNE) at
    ([xshift=-\BookThinkBracketInset,yshift=-#2]frame.north east);
  \coordinate (BookBracketSW) at
    ([xshift=\BookThinkBracketInset,yshift=#3]frame.south west);
  \coordinate (BookBracketSE) at
    ([xshift=-\BookThinkBracketInset,yshift=#3]frame.south east);
  \coordinate (BookBracketLeftTopTip) at (BookBracketNW);
  \coordinate (BookBracketRightTopTip) at (BookBracketNE);
  \coordinate (BookBracketLeftBottomTip) at (BookBracketSW);
  \coordinate (BookBracketRightBottomTip) at (BookBracketSE);
  \ifdim#2>0pt\relax
    \coordinate (BookBracketLeftTopTip) at
      ([xshift=\BookThinkBracketArmLength]BookBracketNW);
    \coordinate (BookBracketRightTopTip) at
      ([xshift=-\BookThinkBracketArmLength]BookBracketNE);
  \fi
  \ifdim#3>0pt\relax
    \coordinate (BookBracketLeftBottomTip) at
      ([xshift=\BookThinkBracketArmLength]BookBracketSW);
    \coordinate (BookBracketRightBottomTip) at
      ([xshift=-\BookThinkBracketArmLength]BookBracketSE);
  \fi
  \draw[draw=#1,line width=\BookThinkBracketRuleWidth,line cap=butt,line join=miter]
    (BookBracketLeftTopTip)--(BookBracketNW)--(BookBracketSW)--(BookBracketLeftBottomTip);
  \draw[draw=#1,line width=\BookThinkBracketRuleWidth,line cap=butt,line join=miter]
    (BookBracketRightTopTip)--(BookBracketNE)--(BookBracketSE)--(BookBracketRightBottomTip);
}

% Η διακόσμηση επιλέγεται μετά τον κοινό renderer. Το interior hidden αφήνει
% να φαίνεται το πραγματικό χαρτί, ακόμη και σε μη λευκό φόντο σελίδας.
% Τα overlays του renderer κρατούν το εικονίδιο μόνο στο πρώτο τμήμα.
\tcbset{
  BookCalloutDecoration/AttentionBox/.style={},
  BookCalloutDecoration/HistoryBox/.style={},
  BookCalloutDecoration/ThinkBox/.style={%
    no borderline,interior hidden,
    underlay unbroken={\BookCalloutBrackets{BookCalloutAccent}{\BookThinkBracketInset}{\BookThinkBracketInset}},
    underlay first={\BookCalloutBrackets{BookCalloutAccent}{\BookThinkBracketInset}{0pt}},
    underlay middle={\BookCalloutBrackets{BookCalloutAccent}{0pt}{0pt}},
    underlay last={\BookCalloutBrackets{BookCalloutAccent}{0pt}{\BookThinkBracketInset}}
  }
}

% Λεπτά περιβλήματα με την παλιά ακριβώς δημόσια σύνταξη.
% Η μόνη διαφορά του δεύτερου είναι το προαιρετικό θέμα.
\newcommand{\BookDeclareCallout}[5]{%
  \newtcolorbox{#1}{BookCalloutRender={#2}{#3}{#4}{#5}{},BookCalloutDecoration/#1}%
}
\newcommand{\BookDeclareTitledCallout}[5]{%
  \newtcolorbox{#1}[1][]{BookCalloutRender={#2}{#3}{#4}{#5}{##1},BookCalloutDecoration/#1}%
}

\BookDeclareCallout{AttentionBox}{BookAlertBackground}{BookCalloutAccent}{\BookCalloutWarning}{\LabelAlert}
\BookDeclareCallout{ThinkBox}{white}{BookCalloutAccent}{\BookCalloutBulb}{\LabelThink}
\BookDeclareTitledCallout{HistoryBox}{BookHistorySoft}{BookCalloutAccent}{\BookCalloutHourglass}{\LabelHistory}
